%!TEX TS-program = lualatex
%!TEX encoding = UTF-8 Unicode

\documentclass[11pt, a4paper]{awesome-cv}

% Configure page margins with geometry
\geometry{left=1.4cm, top=.8cm, right=1.4cm, bottom=1.8cm, footskip=.5cm}

% Specify the location of the included fonts
\fontdir[fonts/]

% Color for highlights
\colorlet{awesome}{awesome-red}
% Uncomment if you would like to specify your own color
% \definecolor{awesome}{HTML}{CA63A8}

% Colors for text
% Uncomment if you would like to specify your own color
% \definecolor{darktext}{HTML}{414141}
% \definecolor{text}{HTML}{333333}
% \definecolor{graytext}{HTML}{5D5D5D}
% \definecolor{lighttext}{HTML}{999999}
% \definecolor{sectiondivider}{HTML}{5D5D5D}

% Set false if you don't want to highlight section with awesome color
\setbool{acvSectionColorHighlight}{true}

% If you would like to change the social information separator from a pipe (|) to something else
\renewcommand{\acvHeaderSocialSep}{\quad\textbar\quad}


%-------------------------------------------------------------------------------
%	PERSONAL INFORMATION
%	Comment any of the lines below if they are not required
%-------------------------------------------------------------------------------
% Available options: circle|rectangle,edge/noedge,left/right
% \photo[rectangle,edge,right]{./examples/profile}
\name{Yara}{Marín}
\position{Desarrolladora Backend {\enskip\cdotp\enskip} Entusiasta de Ciberseguridad}
\address{España}

\email{profesional@miguvt.com}
\homepage{miguvt.com}
\github{MiguVT}
\extrainfo{\faBlog\enskip \href{https://blog.miguvt.com}{blog.miguvt.com}}
% \linkedin{your-linkedin}
% \gitlab{gitlab-id}
% \stackoverflow{SO-id}{SO-name}
% \twitter{@twit}
% \skype{skype-id}
% \reddit{reddit-id}
% \medium{madium-id}
% \kaggle{kaggle-id}
% \hackerrank{hackerrank-id}
% \googlescholar{googlescholar-id}{name-to-display}
%% \firstname and \lastname will be used
% \googlescholar{googlescholar-id}{}
% \extrainfo{extra information}

\quote{``Centrado en backend. Privacidad ante todo. Siempre innovando."}


%-------------------------------------------------------------------------------
\begin{document}

% Print the header with above personal information
% Give optional argument to change alignment(C: center, L: left, R: right)
\makecvheader[C]

% Print the footer with 3 arguments(<left>, <center>, <right>)
% Leave any of these blank if they are not needed
\makecvfooter
  {\today}
  {Yara Marín~~~·~~~Curriculum Vitae}
  {\thepage}


%-------------------------------------------------------------------------------
%	CV/RESUME CONTENT
%-------------------------------------------------------------------------------

%-------------------------------------------------------------------------------
%	RESUMEN
%-------------------------------------------------------------------------------
\cvsection{Resumen}
\begin{cvparagraph}
Desarrolladora backend y estudiante de redes/sistemas dedicada a construir sistemas fiables, escalables y con enfoque en privacidad. Experiencia administrando infraestructuras autoalojadas (Linux, Docker, Nginx) y contribuyendo activamente a comunidades open source. Combino resolución de problemas y automatización con comunicación clara, uniendo arquitectura técnica y creación digital.
\end{cvparagraph}


%-------------------------------------------------------------------------------
%	EXPERIENCIA
%-------------------------------------------------------------------------------
\cvsection{Experiencia}
\begin{cventries}

  \cventry
    {Desarrolladora Web (Autónoma)} % Job title
    {Subastasac} % Organization
    {Remoto} % Location
    {Actualidad} % Date(s)
    {
      \begin{cvitems} % Description(s) of tasks/responsibilities
        \item {Lidero la remodelación completa y modernización del portal web de Subastasac, optimizándolo para subastas online concurrentes.}
        \item {Mejoro la seguridad y la experiencia de usuario para garantizar un sistema de pujas fiable y de alto rendimiento.}
      \end{cvitems}
    }

  \cventry
    {Prácticas de IT (Prácticas de empresa)} % Job title
    {SSHTeam} % Organization
    {España} % Location
    {2025} % Date(s)
    {
      \begin{cvitems} % Description(s) of tasks/responsibilities
        \item {Soporte técnico, diagnóstico de incidencias y colaboración en tareas de desarrollo web durante dos meses de prácticas.}
      \end{cvitems}
    }

\end{cventries}

%-------------------------------------------------------------------------------
%	CONTRIBUCIONES OPEN SOURCE
%-------------------------------------------------------------------------------
\cvsection{Contribuciones Open Source}
\begin{cventries}

  \cventry
    {Contribuidora Principal}
    {\href{https://github.com/Chatterino/chatterino2}{Chatterino2} (2,400+ GitHub Stars)}
    {GitHub}
    {2023 - Actualidad}
    {
      \begin{cvitems}
        \item {Desarrollé e integré soporte de extension ID para plugins personalizados en este cliente de chat de Twitch de alto rendimiento escrito en C++.}
      \end{cvitems}
    }

  \cventry
    {Desarrolladora Open Source}
    {AxMinions (30,000+ Downloads) \& ServerMappings}
    {GitHub}
    {2023 - 2024}
    {
      \begin{cvitems}
        \item {Arreglé problemas críticos de dependencias y resolví bugs técnicos en AxMinions, mejorando la estabilidad para miles de usuarios.}
        \item {Añadí información optimizada de servidores de Minecraft al registro ServerMappings de LunarClient.}
      \end{cvitems}
    }

  \cventry
    {Contribuidora del Ecosistema Linux/VR}
    {\href{https://wiki.vronlinux.org/}{LVRA (Linux VR Adventures) Wiki}}
    {GitHub / Discord}
    {2024}
    {
      \begin{cvitems}
        \item {Escribí un script crítico para solucionar el bug de "blank EDID" en Valve Index bajo Linux, adoptado ampliamente y convertido en estándar dentro de la comunidad (wiki.vronlinux.org).}
      \end{cvitems}
    }
  
  \cventry
    {Contribuidora Activa}
    {FreesmLauncher (380+ GitHub Stars)}
    {GitHub / Discord}
    {2024 - Actualidad}
    {
      \begin{cvitems}
        \item {Contribuyo al código y ayudo a usuarios con troubleshooting técnico dentro del servidor de Discord de la comunidad.}
      \end{cvitems}
    }

\end{cventries}


%-------------------------------------------------------------------------------
%	PROYECTOS
%-------------------------------------------------------------------------------
\cvsection{Proyectos Destacados}
\begin{cventries}

  \cventry
    {Desarrolladora Principal} % Role
    {\href{https://migurinth.miguvt.com/}{Migurinth}} % Title
    {} % Location
    {En curso} % Date(s)
    {
      \begin{cvitems} % Description(s)
        \item {Diseñé un fork de Modrinth orientado a privacidad, eliminando telemetría y anuncios y priorizando la protección de datos del usuario.}
        \item {Implementé integración completa offline adaptada para mayor accesibilidad en servidores.}
      \end{cvitems}
    }

  \cventry
    {Creadora \& Desarrolladora} % Role
    {\href{https://github.com/MiguVT/ChatterinoWatch}{ChatterinoWatch}} % Title
    {} % Location
    {En curso} % Date(s)
    {
      \begin{cvitems} % Description(s)
        \item {Creé una extensión ligera de navegador (Chrome \& Firefox) que sincroniza automáticamente el stream que ves con instancias nativas de Chatterino.}
        \item {Mantengo una base activa de ~160 usuarios diarios que dependen de la extensión para sincronización continua.}
      \end{cvitems}
    }

  \cventry
    {Desarrolladora de Kernel} % Role
    {\href{https://github.com/MiguVT/NP2_Kernel}{NP2\_Kernel}} % Title
    {} % Location
    {2023 - 2024} % Date(s)
    {
      \begin{cvitems} % Description(s)
        \item {Compilé y mantuve un kernel Android personalizado para Nothing Phone 2 basado en LineageOS.}
        \item {Incorporé parches de rendimiento, integración avanzada de drivers y mejoras de seguridad.}
      \end{cvitems}
    }

\end{cventries}


%-------------------------------------------------------------------------------
%	INFRAESTRUCTURA & AUTOALOJADO
%-------------------------------------------------------------------------------
\cvsection{Infraestructura \& Sistemas}
\begin{cventries}
  \cventry
    {Arquitectura Segura de Infraestructura Autoalojada} % Job title
    {Homelab \& Integración Cloud} % Organization
    {} % Location
    {En curso} % Date(s)
    {
      \begin{cvitems} % Description(s) of tasks/responsibilities
        \item {Diseñé y mantengo una infraestructura digital completa conectando una Raspberry Pi 5 local con un VPS de OVH mediante túneles Wireguard.}
        \item {Desplegué y administro servicios autoalojados: Matrix para comunicación descentralizada, Nextcloud, Vaultwarden y bots propios de Discord (YAGPDB + VOCARD).}
        \item {Monté un entorno completo para hosting de aplicaciones y juegos con Pterodactyl Panel, Wings y extensiones personalizadas para un uso óptimo de recursos.}
        \item {Enruto tráfico de forma segura con Nginx (stream/reverse proxy), aplicando principios zero-trust y aislamiento de red robusto.}
      \end{cvitems}
    }
\end{cventries}


%-------------------------------------------------------------------------------
%	HABILIDADES
%-------------------------------------------------------------------------------
\cvsection{Habilidades}
\begin{cvskills}
  \cvskill
    {Idiomas}
    {Español (Nativo) \enskip\cdotp\enskip Inglés (Fluido escrito, avanzado hablado)}
  \cvskill
    {Backend}
    {Python, Node.js, Go, SQL (PostgreSQL), JavaScript, TypeScript, Lua, PHP, C\#, Java}
  \cvskill
    {Frontend}
    {HTML5, CSS3, React, Vue.js, Vite, Tailwind CSS, Webpack}
  \cvskill
    {Infraestructura}
    {Linux (Debian/DietPi, Arch, Kali), Docker, Nginx, Wireguard VPN, Bash/Shell Scripting, Hardware Diagnostics}
  \cvskill
    {Herramientas \& Flujo de trabajo}
    {Git, VSCodium, Unity, Blender, OBS, DaVinci Resolve, GIMP}
\end{cvskills}


%-------------------------------------------------------------------------------
%	EDUCACIÓN
%-------------------------------------------------------------------------------
\cvsection{Educación}
\begin{cventries}

  \cventry
    {Sistemas Microinformáticos y Redes (SMR)} % Degree
    {FDP Rioja} % Institution
    {España} % Location
    {2024 - 2026} % Date(s)
    {
      \begin{cvitems} % Description(s) bullet points
        \item {Enfoque en administración de sistemas, arquitecturas de red y configuración de hardware.}
      \end{cvitems}
    }

  \cventry
    {Educación Secundaria Obligatoria (ESO)} % Degree
    {IES} % Institution
    {España} % Location
    {Finalizada} % Date(s)
    {}

\end{cventries}


%-------------------------------------------------------------------------------
%	CREACIÓN DE CONTENIDO
%-------------------------------------------------------------------------------
\cvsection{Creación de Contenido \& Comunicación}
\begin{cventries}
  \cventry
    {Creadora de Contenido Digital / VTuber}
    {MiguVT}
    {Twitch / YouTube}
    {En curso}
    {
      \begin{cvitems}
        \item {Construyo y hago crecer una comunidad técnica y de gaming con 1.400+ seguidores en Twitch y 730+ suscriptores en YouTube.}
        \item {Aporto comunicación en directo, ideación de proyectos y gestión de comunidad, junto con habilidades creativas (OBS, edición de vídeo y pipeline de avatares/3D).}
      \end{cvitems}
    }
\end{cventries}

\end{document}
